\subsection{Every longer singleton at two outer edges}\label{proof:R2-C048}

\paragraph{R48.0 Complete bounded statement}

Work with the original finite probability space and the original real skew kernel, pointwise bounded in absolute value by one. \(At\) an internal root, maximal genuine equal-arm pairing leaves precisely one original arm of length \(m \ge  3\). \(At\) the other end of an original inner connector of length \(L \ge  1\), the strict residual terminal list has size at least two. Both original vertices may have arbitrary genuine equal-pair decorations. Append exactly two original outer edges from the internal root to the receiving vertex.

Theorem \(T20\). The complete integer direction table in \mbox{\hyperref[part:R48.2]{part R48.2}} gives an actual original message \(H\) at the internal root, and \(U = T^2 H\) at the receiver, such that

\fitdisplay{
\boxed{I_q(U)^3\le \exp(-d/20)}.
}{R48.0.1}

The directions depend only on the original terminal integers, \(m\) and \(L\). The same table applies across the \(d\) domains and spectral frequencies. Every finite terminal cardinality, every finite arm length, every \(L\), every original equal pair and the boundary \(d = 0\) is included.

Alongside the \(m = 2\) case, the \(m \ge  3\) proof uses same-direction ordinary fourths, finite signed coefficients and the ordinary-fourth two-edge bridge with its two \(T\) steps.

The following new ingredients complete the finite catalogue:

\begin{enumerate}
\item The previously exceptional actual \(W_{3} T^2(q f_b)\) direction receives its own fourth estimate, with the true constant-vector cross retained.
\item The actual opposite first-edge family \(\bar f_m A T^{L-1}(q f_{2}),\allowbreak \allowbreak  m > L+1\), receives complete \(L = 2,\allowbreak \allowbreak  L = 3\) and every \(L \ge  4\) fourth estimates.
\item The weak common fourths in C040 are strengthened where needed; its original directions remain unchanged in every analytic domain.
\item For the nonunit odd O terminal, the \(m = 3\), long \(-L\) row retains the complete finite convolution and its last two negative coefficients.
\item The C043 high-domain unit-minimum even directions retain their actual full tail cap and use the correct actual \(U\) fourth, including \(m=3,\allowbreak \allowbreak L=2\).
\end{enumerate}

Theorem (R48.0.1) is an original half-tree interface.

\paragraph{R48.1 Original calculus and the paid two-edge bridge}

\paragraph{R48.1.1 Notation, finite identities and actual positivity}

On the original \(L^2(\mu)\), put

\fitdisplay{
\begin{gathered}
B=-K,\quad u=B\one,\quad Jv=(Ev)\one,\quad T=J+B,\quad A=J-B=T^*,\\
\qquad q=1+u,\quad p=1-u,
\end{gathered}
}{R48.L230}

\fitdisplay{
\begin{gathered}
f_n=A^n\one,\quad \bar f_n=T^n\one,\quad h_n=Tf_{n-1},\quad W_n=A\bar f_{n-1},\\
\qquad \delta_j=\|B^j\one\|^2,\quad d=\delta_1,
\end{gathered}
}{R48.L234}

\fitdisplay{
\begin{gathered}
\theta=\frac{13}{32},\qquad \kappa=\frac{19}{32},\qquad \tau=d-\delta_2,\\
\qquad I_q(V)=E[q^{-\frac{1}{3}}\lvert V\rvert ^{8/3}].
\end{gathered}
}{R48.1.1}

Here \(f_0=\bar f_0=1\), and \(h_n,\allowbreak \allowbreak W_n\) are defined for \(n\ge 1\). In finite propagated sums \(W_0=1\) denotes the genuine zero-edge path.

The original Q3-C014 and Q4-C143 finite calculus gives

\fitdisplay{
\begin{gathered}
0\le d\le1/3,\quad d^2\le \delta_2\le\theta d,\quad\\
\delta_{j+1}\le\theta \delta_j,\quad \|B\|^2\le\theta,\\
\qquad \lvert Bv(x)\rvert \le\|v\|_2.
\end{gathered}
}{R48.1.2}

All A and \(T\) kernels are original nonnegative kernels. All actual finite tree messages are positive on the positive-mass support. Indeed the all-equal labeling at any retained root atom has positive mass and all its diagonal edge factors equal one. In particular \(p,\allowbreak \allowbreak q > 0\) there; the bounds retain the atomwise dependence of \(p\) and \(q\).

Let

\fitdisplay{
\begin{gathered}
G_n(z)=\sum_{i=0}^{n-1}z^i,\quad G_0=0,\qquad\\
\beta_1=\beta_2=0,\quad\\
\beta_n=\sum_{i=1}^{\lfloor(n-1)/2\rfloor}(-1)^{i-1}\delta_i.
\end{gathered}
}{R48.1.3}

Thus \(\beta_{3} = \beta_{4} = d,\allowbreak \allowbreak  \beta_{5} = \beta_{6} = \tau\), and \(\tau \le  \beta_n \le  d\) for \(n \ge  3\). These \(\beta\) values are the actual finite alternating moments; they are not \(1 - E \bar f_{n-1}\). The original finite positive convolution is

\fitdisplay{
1+G_n(B)u=\bar f_n+\sum_{r=3}^n\beta_r \bar f_{n-r}.
}{R48.1.4}

Its completely reversed form holds as well. Every removed term in (R48.1.4) is a nonnegative actual path term. We use this finite identity below.

On the original spectral measure \(\nu\) of \(u\), write \(z = i \sqrt{x},\allowbreak \allowbreak  0 \le  x \le  \theta\). Its mass is \(d\) and \(\int  x^i d\nu = \delta_{i+1}\). For a real polynomial \(P = E_P(x) + z C_P(x)\),

\fitdisplay{
E[P(B)u]=-\int C_P\,d\nu,\qquad
 E[BP(B)u]=-\int E_P\,d\nu.
}{R48.1.5}

Write \(G_n = E_n + z C_n\). For \(n \ge  2\) both coordinates lie in \([1-x,\allowbreak \allowbreak 1]\); for \(n \ge  3,\allowbreak \allowbreak  E_n \le  h(x) := 1-x+x^2\). Also

\fitdisplay{
\|G_n(B)\|\le1\quad(n=0,1\text{ or }n\ge3),\qquad
 \|G_2(B)\|^2\le1+\theta.
}{R48.1.6}

 Lemma~\ref{lem:outer-unit:fractional-tail},
\eqref{eq:outer-unit:deleted-z}, proves \(\lvert G_n(z)-z\rvert\le1\)
for every finite \(n\ge2\), including the short indices.

For pure paths \(\mu_n := E f_n = E \bar f_n \le  1\), with \(\mu_{2} = 1-d,\allowbreak \allowbreak  \mu_{3} = 1-2d\), and \(\mu_n \le  1-2d\) for every \(n \ge  3\). Every \(f_n,\allowbreak \allowbreak \bar f_n,\allowbreak \allowbreak h_n,\allowbreak \allowbreak W_n\) lies in [0,2], and

\fitdisplay{
f_n+h_n=\bar f_n+W_n=2\mu_{n-1},\qquad
 f_nh_n\le1,\quad \bar f_nW_n\le1.
}{R48.1.7}

The means of \(h_{2}\) and \(W_{2}\) are \(1+d\); use their exact low-index formulas.

The all-real signed exponential input in the supplied finite calculus is

\fitdisplay{
E e^{4BV}\le\exp\{4E(BV)+(\frac{17}{4})\|V\|^2\}.
}{R48.1.8}

It applies to every real \(V\) with either sign of the actual mean, including a genuine constant component. 

If actual nonnegative factors \(X,\allowbreak \allowbreak Y\) have \(X+Y \le  2+c+BV\), AM-GM therefore gives

\fitdisplay{
\log E(XY)^4\le4(c+E BV)+(\frac{17}{4})\|V\|^2.
}{R48.1.9}

The upper base is nonnegative because it bounds an actual nonnegative sum. If the second factor has a prefactor \(\exp(-\Gamma)\), add \(-4 \Gamma\); if its prefactor is \(a_0^2\), add \(8 \log a_0\). This fixes the fourth convention used throughout the manuscript. Every expression such as \(E[X]^4\) below denotes the expectation \(E(X^4)\), never the fourth power of the mean.

\paragraph{R48.1.2 The exact original two-edge transfer}

The actual C094, Section 2, gives for every real original \(H\)

\fitdisplay{
\begin{aligned}
 I_q(T^2H)^3
 &\le [1-2d^2+(8-2d)\delta_2]
       (1+1024\delta_2/855)^8(E\lvert H\rvert ^4)^2\\
 &\le \exp(c_2d-5d^2)(E\lvert H\rvert ^4)^2,
 \qquad c_2=\frac{24427}{3420}<36/5.                         
 \end{aligned}
}{R48.1.10}

For clarity the positive-operator step is exact. If \(P\) is positive with row \(r = P\one\), row Cauchy and \(H\)ölder give

\fitdisplay{
\begin{gathered}
I_q(PH)^3\le E(r^4/q)\{E[P(H^2)^2]\}^2\\
\le E(r^4/q)\|P\|^4(E\lvert H\rvert ^4)^2.
\end{gathered}
}{R48.1.11}

Here \(P(H^2)^2\) denotes the square of \(P(H^2)\), not \(P(H^4)\). For \(P=T^2\) the actual row is \(\bar f_{2}=q+Bu\). C094's correlated cubic row calculation gives \(E(\bar f_{2}^4/q) \le  1-2d^2+(8-2d)\delta_2\). Its actual ordinary norm is \(\lVert T^2\rVert  \le  (1+1024 \delta_2/855)^2\). Hence the power in (R48.1.10) is eight. The variance \(\delta_2 = \lVert B^2\one\rVert ^2\) has not been changed to \(d^2\).

For the curvature, the same calculation proves on the whole interval \(0 \le  d \le  \frac{1}{3}\)

\fitdisplay{
\log(1+13d/4-45d^2/16)
 \le13d/4-(\frac{11757}{2720})d^2,
}{R48.L370}

\fitdisplay{
8\log(1+416d/855)\le(\frac{3328}{855})d-3d^2/4.
}{R48.1.12}

Their linear sum is \(c_2\) and their quadratic sum exceeds five. The complete inverse-fourth row and the ordinary norm occur in the same original measure.

In particular every same-direction fourth bound \(E H^4 \le  \exp(-29d/8)\) implies (R48.0.1), since

\fitdisplay{
2(\frac{29}{8})-\frac{36}{5}=1/20.
}{R48.1.13}

Some rows below instead use their explicitly retained negative curvature.

\paragraph{R48.2 The complete structural direction table}\label{part:R48.2}

The following defines the actual unpaired \(H\). Genuine equal pairs are restored at their original vertices by the common restoration argument. A bar means reversal of every arm in that displayed terminal configuration. This specifies the terminal reversal in the direction table.

\paragraph{R48.2.1 Even terminal residuals}\label{part:R48.2.1}

Write an even strict terminal as \(a_{1} < b_{1} < \ldots < a_{j} < b_{j}\).


\begingroup\small
\setlength{\tabcolsep}{3pt}
\setlength{\LTleft}{0pt}\setlength{\LTright}{\fill}
\begin{longtable}{@{}p{0.2200\linewidth}p{0.7000\linewidth}@{}}
\caption{Every longer singleton at two outer edges: parameter values.}\label{tab:parameters:r2-c048:1}\\
\toprule
Entry & Actual \(H\) \\
\midrule\endfirsthead
\toprule
Entry & Actual \(H\) \\
\midrule\endhead
\bottomrule\endfoot
Exact \((1,\allowbreak \allowbreak b),\allowbreak \allowbreak  L=1\) & \(f_m T(p \bar f_b)\). \\
\addlinespace[2pt]
Exact \((1,\allowbreak \allowbreak b),\allowbreak \allowbreak  m=3,\allowbreak \allowbreak L=2\) & \(f_{3} T A(p \bar f_b)\). \\
\addlinespace[2pt]
Exact \((1,\allowbreak \allowbreak b)\), all remaining \(L\ge 2\) & \(W_m T^L(q f_b)\). \\
\addlinespace[2pt]
Exact \((2,\allowbreak \allowbreak b),\allowbreak \allowbreak  b\ge 3,\allowbreak \allowbreak L=1\) & \(f_m T(f_{2} \bar f_b)\). \\
\addlinespace[2pt]
Exact \((2,\allowbreak \allowbreak b),\allowbreak \allowbreak  b\ge 3,\allowbreak \allowbreak L\ge 2\) & \(W_m T^L(\bar f_{2} W_b)\). \\
\addlinespace[2pt]
Exact \((3,\allowbreak \allowbreak b),\allowbreak \allowbreak  b\ge 4,\allowbreak \allowbreak L=1\) & \(f_m T(f_{3} h_b)\). \\
\addlinespace[2pt]
Exact \((3,\allowbreak \allowbreak b),\allowbreak \allowbreak  b\ge 4,\allowbreak \allowbreak L\ge 2\) & \(W_m T^L(f_{3} h_b)\). \\
\addlinespace[2pt]
\(a_{1}\ge 4\), every \(j\ge 1; m\le L\) & \(W_m T^L(F),\allowbreak \allowbreak  F=\prod h_{a_{i}} f_{b_{i}}\). \\
\addlinespace[2pt]
\(a_{1}\ge 4\), every \(j\ge 1; m>L\) & \(\bar f_m A T^{L-1}(F)\), same \(F\). \\
\addlinespace[2pt]
\(a_{1}=2,\allowbreak \allowbreak j\ge 2,\allowbreak \allowbreak L=1\) & \(\bar f_m A(F),\allowbreak \allowbreak  F=\prod h_{a_{i}} f_{b_{i}}\). \\
\addlinespace[2pt]
\(a_{1}=2,\allowbreak \allowbreak j\ge 2,\allowbreak \allowbreak L\ge 2\) & \(W_m T^L(F)\), same \(F\). \\
\addlinespace[2pt]
\(a_{1}=3,\allowbreak \allowbreak j\ge 2,\allowbreak \allowbreak m\le L\) & \(W_m T^L(F),\allowbreak \allowbreak  F=\prod h_{a_{i}} f_{b_{i}}\). \\
\addlinespace[2pt]
\(a_{1}=3,\allowbreak \allowbreak j\ge 2,\allowbreak \allowbreak m>L\) & \(\bar f_m A T^{L-1}(F)\), same \(F\). \\
\addlinespace[2pt]
\(a_{1}=1,\allowbreak \allowbreak j\ge 2,\allowbreak \allowbreak L=1\) & \(f_m T(\overline F),\allowbreak \allowbreak F=q f_{b_{1}}\prod_{i\ge2}h_{a_{i}}f_{b_{i}}\). \\
\addlinespace[2pt]
\(a_{1}=1,\allowbreak \allowbreak j\ge 2,\allowbreak \allowbreak L\ge 2\) & \(W_m T^L(F)\), same \(F\). \\
\end{longtable}
\endgroup


The exact \(-(1,\allowbreak \allowbreak b) m=3,\allowbreak \allowbreak L=2\) turn is the old Q5-C087 short-outer direction. It is used here because the outer length is two. The unit-minimum multipair \(m=3,\allowbreak \allowbreak L=2\) row instead keeps its actual C043 all \(-T\) direction; \mbox{\hyperref[part:R48.3]{part R48.3}} pays that different message explicitly.

\paragraph{R48.2.2 Odd terminal residuals of nonunit minimum}

Write the strict list as \((\ell,\allowbreak \allowbreak r_1,\allowbreak \allowbreak t_{1},\allowbreak \allowbreak \ldots,\allowbreak \allowbreak r_{j},\allowbreak \allowbreak t_{j})\), where \(2 \le  \ell < r_1 < t_{1} < \ldots < r_{j} < t_{j}\) and \(j \ge  1\). Put

\fitdisplay{
F=\begin{cases}
 \bar f_\ell\prod_{i=1}^j h_{r_i}f_{t_i},&r_1\ge4,\\
 \bar f_2f_3h_{t_1}\prod_{i=2}^j h_{r_i}f_{t_i},&r_1=3.
 \end{cases}
}{R48.2.1}

The second row necessarily has \(\ell=2\). Set \(n=\ell+L\) and use

\fitdisplay{
H=\begin{cases}
 W_mT^L F,&m\le n,\\
 \bar f_mAT^{L-1}F,&m>n.
 \end{cases}
}{R48.2.2}

These are exactly the R2-C030 actual O directions.

\paragraph{R48.2.3 Odd terminal residuals of unit minimum}

Write the strict list as \((1,\allowbreak \allowbreak r_1,\allowbreak \allowbreak t_{1},\allowbreak \allowbreak \ldots,\allowbreak \allowbreak r_{j},\allowbreak \allowbreak t_{j}),\allowbreak \allowbreak  j \ge  1\). The actual terminal is the C040 terminal:

\fitdisplay{
F=\begin{cases}
 q\prod_{i=1}^j h_{r_i}f_{t_i},&r_1\ge4,\\
 q f_3h_{t_1}\prod_{i=2}^j h_{r_i}f_{t_i},&r_1=3,\\
 q f_2h_eh_3f_4,&\text{exact list }(1,2,3,4,e),\ j=2,\\
 q f_2h_{t_1}\prod_{i=2}^j h_{r_i}f_{t_i},&\text{all other }r_1=2.
 \end{cases}
}{R48.2.3}

For \(L=1\) use \(H=W_m TF\) when \(r_1\ge 3\), and \(H=\bar f_m AF\) when \(r_1=2\). For every \(L\ge 2\) use

\fitdisplay{
H=\begin{cases}
 W_mT^LF,&m\le L+1,\\
 \bar f_mAT^{L-1}F,&m>L+1.
 \end{cases}
}{R48.2.4}

The crossing in (R48.2.3) is only the exact five-arm list. The same prefix inside a list with at least seven residual arms stays canonical, exactly as in C040. The equality \(m=L+1\) belongs to the first row of (R48.2.4). Each analytic split below estimates the direction fixed by this table.

In every \(A T^{L-1}\) row the receiver-adjacent original connector edge is A and the other \(L-1\) edges are \(T\). Every \(f,\allowbreak \allowbreak \bar f,\allowbreak \allowbreak h,\allowbreak \allowbreak W\) is its actual path of the displayed original length. The special \(T\) A row contains the two original inner edges and their actual common midpoint. Finally every row appends the same two original outer \(T\) edges, giving \(U=T^2H\).

\paragraph{R48.3 The actual \(W_{3}\label{part:R48.3} T^2(q f_b)\) boundary}

The turn \(f_{3} T A(p \bar f_b)\) differs from the C043 all \(-T\) message estimated here.

\paragraph{R48.3.1 Terminal two: the real constant cross and a quadratic reserve}

The actual identities are \(f_{2}=p+Bu\) and \(q+f_{2}=2+Bu\). Therefore

\fitdisplay{
\begin{gathered}
q f_2\le(1+Bu/2)^2,\\
\qquad T(q f_2)\le q-\gamma+B^2u,\\
\qquad\gamma=d-\delta_2/2.
\end{gathered}
}{R48.3.1}

Indeed \(T(Bu)=-d+B^2u\) and \(T((Bu)^2)\le 2\delta_2\). Applying the actual positive \(T\) once more gives

\fitdisplay{
T^2(q f_2)\le \bar f_2-\gamma q+B^3u.
}{R48.3.2}

The exact \(W_{3}\) is \(1-d-u-Bu-B^2u\), so \(W_{3}+\bar f_{2}=2-d-B^2u\). For the sum of \(W_{3}\) and the nonnegative upper in (R48.3.2), use

\fitdisplay{
c=-d-\gamma,\qquad V=-\gamma\,1-Bu+B^2u,
 \qquad E(BV)=\delta_2.
}{R48.3.3}

The vector is not centered. Its genuine constant cross is

\fitdisplay{
\begin{aligned}
 \|V\|^2
 &=\delta_2+\delta_3+\gamma^2-2\gamma d\\
 &=\delta_2+\delta_3-d^2+\delta_2^2/4.                             
 \end{aligned}
}{R48.3.4}

Equation (R48.1.9) gives

\fitdisplay{
\begin{gathered}
\log E[W_3T^2(q f_2)]^4\\
\le-8d+(\frac{41}{4})\delta_2+(\frac{17}{4})\delta_3-(\frac{17}{4})d^2+(\frac{17}{16})\delta_2^2
\end{gathered}
}{R48.L527}

\fitdisplay{
\begin{gathered}
\le-\frac{12839}{4096}d\\
-\frac{17}{4}(1-\theta^2/4)d^2\\
\le-\frac{25}{8}d-4d^2.
\end{gathered}
}{R48.3.5}

The last comparisons are exact: \(12839>12800\), and \(\theta^2<\frac{1}{6}\) gives \((\frac{17}{4})(1-\theta^2/4)>\frac{391}{96}>4\). Thus every original host is included, not only the high domain in which this estimate will later be used.

\paragraph{R48.3.2 Every \(b \ge  3\): one finite terminal polynomial}

Set

\fitdisplay{
\begin{gathered}
R=B G_{b-1}(-B)u,\quad a=1-\beta_b/2,\quad k=a^{-1},\\
\quad V_R=ER^2,\quad t_0=\beta_{b+1},
\end{gathered}
}{R48.L546}

\fitdisplay{
e=d-\beta_{b+2},\quad \rho=E[(Bu)R],
 \qquad\gamma=kt_0-k^2V_R/2.
}{R48.3.6}

The actual finite positive terminal square, also used in C087 parts R48.3–R48.4, is \(q f_b \le  a^2(1+kR/2)^2\). Its first and second rows give

\fitdisplay{
\begin{gathered}
T(q f_b)\le a^2[q-\gamma+kBR],\\
\qquad T^2(q f_b)\le a^2[\bar f_2-\gamma q-ke+kB^2R].
\end{gathered}
}{R48.3.7}

Both displayed upper factors come from actual positive kernels. The exact same-parameter bounds in C087 (4.3), at every finite \(b\), are

\fitdisplay{
\begin{gathered}
1\le k\le6/5,\quad 3d/10\le\gamma\le6d/5,\\
\quad0\le ke\le39d/80,\quad \beta_b\ge\tau.
\end{gathered}
}{R48.3.8}

We use simultaneous consequences of the definitions in (R48.3.6). For \(b\ge 4\) the finite \(G\) contraction gives \(\lVert BR\rVert ^2\le \delta_3\). For \(b=3\) the actual \(G_{2}\) gives \(\lVert BR\rVert ^2=\delta_3+\delta_4\).

Write \(G_{b-1}(-z)=E-zC\), so \(R=(xC+zE)u\). Then

\fitdisplay{
\begin{gathered}
e=\int xC\,d\nu,\quad \rho=\int xE\,d\nu,\\
\quad \langle Bu,BR\rangle=\int x^2C\,d\nu\ge0.
\end{gathered}
}{R48.3.9}

The entire finite difference \(\rho-e\) is controlled without a limiting series. If \(b-1\) is even, \(E=C\). If \(b-1=2r+1\) is odd, \(E-C=(-x)^r\); here \(r\ge 1\). A positive difference requires even \(r\ge 2\), whence

\fitdisplay{
\rho-e\le \delta_4.
}{R48.3.10}

For \(b=3\) specifically \(\rho=e=\delta_2\).

Apply AM-GM to \(W_{3}\) and the bracket in the second row of (R48.3.7). The scalar and vector are

\fitdisplay{
c=-d-\gamma-ke,\qquad V=-\gamma\,1-Bu+kBR,
 \quad c+EBV=-d-\gamma+k(\rho-e).
}{R48.3.11}

For \(V_0=-Bu+kBR\), its true mean is \(EV_{0}=d-\) ke \(\ge 41d/80\). Consequently its constant correction has the whole-domain bound

\fitdisplay{
\begin{gathered}
\gamma^2-2\gamma EV_0\\
\le\max_{0\le y\le6d/5}\{y^2-(\frac{41}{40})dy\}\\
\le21d^2/100\le7d/100.
\end{gathered}
}{R48.3.12}

The maximum is at an endpoint of this convex quadratic. \(At\) the nonzero endpoint its coefficient is \(\frac{36}{25}-\frac{123}{100}=\frac{21}{100}\). Using the nonnegative cross in (R48.3.9), for \(b\ge 4\) we obtain

\fitdisplay{
\begin{gathered}
\|V\|^2\le \delta_2+(\frac{36}{25})\delta_3+7d/100\\
\le(\frac{4569}{6400})d<3d/4.
\end{gathered}
}{R48.3.13}

For \(b=3\) the same argument gives a norm below \(d\): indeed \(\theta^2(1+\theta)<\frac{1}{4}\) and \(\theta+(\frac{36}{25})(\frac{1}{4})+\frac{7}{100}<1\).

Now \(8 \log a \le  -4 \beta_b\). For \(b\ge 4\), equations (R48.1.9) and (R48.3.8)-(R48.3.13) give a fourth loss strictly larger than

\fitdisplay{
\begin{gathered}
4\kappa+4+\frac{6}{5}-(\frac{24}{5})\theta^3-\frac{51}{16}\\
>\frac{19}{8}+\frac{26}{5}-\frac{3}{8}-\frac{51}{16}=\frac{321}{80}>4.
\end{gathered}
}{R48.3.14}

Here \(\theta^3<\frac{5}{64}\) pays the penultimate strict comparison. \(At b=3,\allowbreak \allowbreak  \beta_b=d\) and \(\rho=e\); the weaker norm bound \(d\) gives loss at least \(4+4+\frac{6}{5}-\frac{17}{4}=\frac{99}{20}>4\). We have therefore proved

\fitdisplay{
\boxed{E[W_3T^2(q f_b)]^4\le e^{-4d}\quad(b\ge3).}
}{R48.3.15}

Every finite \(b\), both parity classes, the true \(G_{2}\) exception and the actual constant cross are included in this unchanged all \(-T\) direction.

\paragraph{R48.4 The actual \(\bar f_m A T^{L-1}\label{part:R48.4}(q f_{2})\) family}

Assume \(L\ge 2\) and \(m>L+1\), precisely the second branch of (R48.2.4). Write \(J_{m,\allowbreak \allowbreak L}=\bar f_m A T^{L-1}(q f_{2})\). We prove

\fitdisplay{
E J_{m,L}^4\le
 \begin{cases}
 e^{-5d},&L=2,\\
 e^{-167d/48},&L=3,\\
 e^{-6d},&L\ge4.
 \end{cases}
}{R48.4.1}

The positive-message calculation uses the stated single-root orientation; the forward \(U\) family is different.

\paragraph{R48.4.1 \(L=2\) and every \(m\ge 4\)}

Apply A to (R48.3.1):

\fitdisplay{
AT(q f_2)\le W_2-\gamma p-B^3u.
}{R48.4.2}

Keep the last two positive coefficients of (R48.1.4 \(): \bar f_m \le  1+G_m u-\beta_m-\beta_{m-1}q\). Together with \(W_{2}=p-Bu\), the sum of \(\bar f_m\) and the upper in (R48.4.2) is at most \(2+c+BV\), where

\fitdisplay{
c=-(\beta_m+\beta_{m-1}+\gamma),\quad
 V=B(G_{m-2}(B)-B)u+(\gamma-\beta_{m-1})1.
}{R48.4.3}

For \(n=m-2\ge 2\) the finite coordinate \(C_n\le 1\) gives

\fitdisplay{
EBV=\int x(C_n-1)\,d\nu\le0.
}{R48.4.4}

Moreover \(\lvert \gamma-\beta_{m-1}\rvert \le \delta_2/2\), because both actual \(\beta\) values lie in \([d-\delta_2,\allowbreak \allowbreak d]\). The finite inequality \(\lvert G_n-z\rvert \le 1\) yields

\fitdisplay{
\|V\|\le\sqrt{\delta_2}+\delta_2/2
 <(\frac{19}{24})\sqrt d.
}{R48.4.5}

For the rational comparison use \(\sqrt{\theta}<\frac{2}{3},\allowbreak \allowbreak  \theta<\frac{5}{12}\) and \(\sqrt{d}<\frac{3}{5}\), so the second normalized term is less than \(\frac{1}{8}\). The scalar obeys \(c\le -[3d-(\frac{5}{2})\delta_2]\le -127d/64\). Equation (R48.1.9) now gives fourth loss at least

\fitdisplay{
\frac{127}{16}-(\frac{17}{4})(\frac{19}{24})^2
 =\frac{12151}{2304}>5.
}{R48.4.6}

The constant in (R48.4.3) was included in the norm, not discarded.

\paragraph{R48.4.2 \(L=3\) and every \(m\ge 5\)}

From (R48.3.1),

\fitdisplay{
AT^2(q f_2)\le W_3-\gamma W_2+\delta_2-B^4u.
}{R48.4.7}

The actual finite unequal-arm sum for \(h=3,\allowbreak \allowbreak  H=m\) is

\fitdisplay{
W_3+\bar f_m\le2-(2d-\delta_2)+B^3G_{m-3}(B)u.
}{R48.4.8}

Set \(n=m-3\ge 2\). Combining (R48.4.7)-(R48.4.8) gives

\fitdisplay{
\begin{gathered}
c=-3d+(\frac{5}{2})\delta_2,\quad\\
V=\gamma\one+V_0,\quad V_0=\gamma u+B^2(G_n(B)-B)u.
\end{gathered}
}{R48.4.9}

The exact finite coordinates of \(V_0\) are \(\gamma-x E_n+z x(1-C_n)\). Thus

\fitdisplay{
EV_0=\int x(C_n-1)\,d\nu\le0,
 \qquad\|V_0\|^2\le\gamma^2d+\delta_3+\delta_6.
}{R48.4.10}

The omitted real cross is \(-2 \gamma \int  x E_n\) and is nonpositive. Since \(0\le \gamma\le d\), the full constant cross is also nonpositive, and

\fitdisplay{
\begin{gathered}
\|V\|^2\le\gamma^2(1+d)+\delta_3+\delta_6\\
\le[\frac{4}{9}+\theta^2+\theta^5]d<2d/3.
\end{gathered}
}{R48.4.11}

One entirely rational last check is \(\theta^2<\frac{1}{6},\allowbreak \allowbreak  \theta^5<\frac{1}{32}\), and \(\frac{4}{9}+\frac{1}{6}+\frac{1}{32}=\frac{185}{288}<\frac{2}{3}\). The actual mean in the exponent satisfies

\fitdisplay{
\begin{gathered}
EBV=-\gamma d+E[B^3(G_n-B)u]\le \delta_2,\\
\quad c+EBV\le-3d+(\frac{7}{2})\delta_2\le-101d/64.
\end{gathered}
}{R48.4.12}

Here \(E[B^3(G_n-B)u]=\int  x E_n\le \delta_2\). The fourth loss is therefore at least

\fitdisplay{
\frac{101}{16}-\frac{17}{6}=167/48.
}{R48.4.13}

In particular no assertion stronger than the proved (R48.4.13) is needed. On \(d\ge \frac{1}{5},\allowbreak \allowbreak  (\) R48.1.10) already gives an \(M=2\) loss at least

\fitdisplay{
2(\frac{167}{48})+1-\frac{36}{5}=\frac{91}{120}>1/20.
}{R48.4.14}

\paragraph{R48.4.3 Every finite \(L\ge 4\)}

Finite Duhamel applied to the first row (R48.3.1) gives exactly

\fitdisplay{
\begin{gathered}
AT^{L-1}(q f_2)\le W_L-\gamma W_{L-1}-B^{L+1}u\\
+\sum_{i=0}^{L-2}(EB^{i+2}u)W_{L-2-i}.
\end{gathered}
}{R48.4.15}

The last term has \(W_{0}=1\). The \(i=0\) coefficient is zero. The positive classes of the remaining finite sum are \(\delta_2,\allowbreak \allowbreak \delta_4,\allowbreak \allowbreak \delta_6\),...; hence, using \(0\le W_j\le 2\), their entire finite contribution is at most

\fitdisplay{
C\le\frac{2\delta_2}{1-\theta^2}\le\frac{832}{855}d<d.
}{R48.4.16}

We may discard the negative \(-\gamma W_{L-1}\). This is after actual positive propagation, not an unsupported sign change of an input. Since \(4\le L<m\), \eqref{eq:r40:remaining-bare-strong} gives a vector \(Z\) with

\fitdisplay{
c_{\rm bare}+E(BZ)\le-219d/70,\qquad
 \|Z\|\le\sqrt d/3.
}{R48.4.17}

The new tail vector is \(-B^L u\). Its norm is at most \(\theta^2 \sqrt{d} <\sqrt{d}/6\) and its exponent mean is at most \(d/6\) in absolute value. Combining it with \(Z\) and (R48.4.16), the complete sum has

\fitdisplay{
c+EBV\le(-\frac{219}{70}+1+\frac{1}{6})d=-206d/105,
 \qquad\|V\|\le\sqrt d/2.
}{R48.4.18}

The fourth loss is at least \(\frac{824}{105}-\frac{17}{16}=\frac{11399}{1680}>6\). This proves the last row of (R48.4.1) for every finite \(L\ge 4\). In particular the last scalar in (R48.4.15), the \(G_{2}\) boundaries in the bare vector and all later residue classes remain.

\paragraph{R48.4.4 High-score consequence in the exact C040 directions}\label{part:R48.4.4}

For \(c\ge 5\) and \(d\ge \frac{1}{5}\), actual path means give

\fitdisplay{
h_c\le2\mu_{c-1}\le2-5d\le1.
}{R48.4.19}

The slightly stronger supplied \(\mu_{c-1}\le 1-83d/32\) for \(c-1\ge 4\) justifies the weaker \(1-5d/2\) used here. For \(c=4\) and \(d\ge \frac{1}{4},\allowbreak \allowbreak  h_{4}\le 2 \mu_{3}=2-4d\le 1\). Thus in those exact domains

\fitdisplay{
q f_2h_c\le q f_2.
}{R48.4.20}

If \(m\le L+1\), positive propagation compares the C040 message to \(W_m T^L(q f_{2})\). Except at \(m=3,\allowbreak \allowbreak L=2\), its actual C087 fourth loss is at least \(\frac{3487}{640}>\frac{29}{8}\). \(At m=3,\allowbreak \allowbreak L=2,\allowbreak \allowbreak  (\) R48.3.5) and (R48.1.10) give, on \(d\ge \frac{1}{5}\),

\fitdisplay{
\log I_q(T^2H)^3
 \le(\frac{36}{5}-\frac{25}{4})d-13d^2
 \le-33d/20.
}{R48.4.21}

If \(m>L+1\), the actual comparator is \(J_{m,\allowbreak \allowbreak L}\), exactly (R48.4.1). Its \(L=2\) and all \(L\ge 4\) fourths pay (R48.1.13 \(); L=3\) is paid by (R48.4.14). Consequently all original \(L\ge 2\), all \(m\ge 3\) and both entire high domains in (R48.4.19)-(R48.4.20) satisfy (R48.0.1) with their unchanged C040 directions.

\paragraph{R48.5 Every unit-minimum odd terminal}\label{part:R48.5}

Use Proposition~\ref{prop:r40:ordinary-catalogue} and the exact-triple calculation in \hyperref[part:outer-unit:triples]{the ordinary triple construction}, with the fixed direction table (R48.2.3)--(R48.2.4).

\paragraph{R48.5.1 First pair three and exactly one pair}
For \(F=qf_3h_c\), \(c\ge4\), the receiver and connector table is
exactly the one in \eqref{eq:r40:remaining-three-single}.
That full-domain same-direction calculation gives
\[
 \E H^4\le e^{-9d/2}\le e^{-39d/10}.
 \tag{R48.5.3}
\]
Genuine equal-pair factors are restored by positivity.
Since \(9/2>29/8\), the ordinary-fourth two-edge bridge
(R48.1.10)--(R48.1.13) proves the required rate \(1/20\) for the
actual \(U=T^2H\). Its \(q^{-1/3}\) denominator is unchanged.

\paragraph{R48.5.2--5.5 Unit-minimum triples with an outer path}
\label{part:outer-unit:outer}
The messages \(F_c=qf_2h_c\), \(c\ge3\), and their actual
\(m,L\) direction table coincide with
\hyperref[part:outer-unit:triples]{the ordinary triple construction}.
Its simultaneous fourth and eighth calculation includes both short
connector directions and every \(m=3\) row; no terminal is reoriented here.

The transfer from a same-direction ordinary fourth
has a formulation with an arbitrary outer length \(r\ge1\).
Set \(P_r=T^r\), \(\rho_r=P_r\one\), with the ordinary \(L^2\)
operator norm in the original measure. Positive-row Cauchy and
H\"older give
\begin{equation}\label{eq:outer-unit:outer-transfer}
\begin{aligned}
 I_q(P_rH)^3
 &\le \E(\rho_r^4/q)\{\E[(P_r(H^2))^2]\}^2\\
 &\le \E(\rho_r^4/q)\|P_r\|^4(\E H^4)^2.
\end{aligned}
\end{equation}
Indeed \(|P_rH|^2\le\rho_rP_r(H^2)\), and the exponents \(3,3/2\)
separate \((\rho_r^4/q)^{1/3}\) and
\((P_r(H^2)^2)^{2/3}\). Thus the formula retains the inverse row
weight at the receiving vertex; it does not turn the ordinary moment
of \(H\) into a weighted eighth moment.
At \(r=2\), the previously established actual row and norm estimates
(R48.1.10) give
\begin{equation}\label{eq:outer-unit:two-edge-test}
 \E H^4\le e^{-\eta d}\quad\Longrightarrow\quad
 \log I_q(T^2H)^3\le(c_2-2\eta)d-5d^2,
 \qquad c_2<36/5.
\end{equation}
The theorem below uses this explicit two-edge value, rather than an
unproved contraction at arbitrary outer lengths.

For \(c=3\), the complete ordinary proof gives fourth rate
\(4667/1280>29/8\) at \(L=1\), and rate greater than four for
every \(L\ge2\). Hence \eqref{eq:outer-unit:two-edge-test} proves
the desired rate \(1/20\) for every \(m,L,d\).
For \(c=4\), the same proof gives rate \(39/10\) at \(L=1\).
At \(L\ge2,d\le1/4\), every row of
\eqref{eq:outer-unit:connector-table} pays \(29/8\), except
\(L\ge4,m\ge4,1/5\le d\le1/4\).
On that remaining interval retain the actual fourth rate
\(2369/700\) from \eqref{eq:outer-unit:curved-long-four}; then
\[
 2(2369/700)+5d-36/5\ge199/350>1/20.
\]
On \(d\le1/5\), the smaller first-row error in
\eqref{eq:outer-unit:four-small-error} instead gives fourth rate
\(7807/2100>37/10\). On \(d\ge1/4,L\ge2\), the actual cap
\(h_4\le1\) permits the same-direction comparison in
\hyperref[part:R48.4.4]{the high-score calculation}, including its
curved \(W_3T^2\) and opposite-first-edge \(L=3\) cases.

\paragraph{Every exact \(12c\) with \(c\ge5\)}\label{part:R48.5.5}
On \(d\le1/5\), the unit-connector rate is \(19/4\) and every
nonunit-connector fourth rate exceeds \(37/10\), by
\eqref{eq:outer-unit:connector-table} and its unit-connector calculation.
On \(1/5\le d\le1/4,L=1\), the single joint calculation
\eqref{eq:outer-unit:middle-joint} gives rate \(321/80>4\).
These rates all pay \eqref{eq:outer-unit:two-edge-test}.
On \(d\ge1/4,L=1\), use the actual pointwise and mean estimate
\eqref{eq:outer-unit:high-unit}: \(\E H^4<5/16\), while
\(c_2d-5d^2<2\). Thus
\begin{equation}\label{eq:outer-unit:high-outer}
 I_q(T^2H)^3<8(5/16)^2=25/32<5/6\le e^{-d/2}.
\end{equation}
Finally, on \(d\ge1/5,L\ge2\), the actual cap \(h_c\le1\)
gives exactly the forward or opposite-first-edge comparator in
\hyperref[part:R48.4.4]{the high-score calculation}; it proves the
same rate \(1/20\).
These intervals meet at their endpoints and retain the prescribed
integer direction throughout. Positivity restores all original equal
pairs before both outer \(T\) steps. At \(d=0\) the message and both
outer rows are one.

\paragraph{R48.5.6 The remaining complete unit-minimum odd catalogue}

The following rows are the same-direction fourth entries of Table~\ref{tab:r40:ordinary-losses}:


\begingroup\small
\setlength{\tabcolsep}{3pt}
\setlength{\LTleft}{0pt}\setlength{\LTright}{\fill}
\begin{longtable}{@{}p{0.2200\linewidth}p{0.7000\linewidth}@{}}
\caption{Every longer singleton at two outer edges: parameter values.}\label{tab:parameters:r2-c048:5}\\
\toprule
Entry & Existing fourth loss \\
\midrule\endfirsthead
\toprule
Entry & Existing fourth loss \\
\midrule\endhead
\bottomrule\endfoot
\(r_1\ge 4,\allowbreak \allowbreak  j\ge 1\) & \(5+(\frac{64}{5})(j-1)\). \\
\addlinespace[2pt]
\(r_1=3,\allowbreak \allowbreak  j\ge 2\) & \(7+(\frac{64}{5})(j-2)\). \\
\addlinespace[2pt]
exact \((1,\allowbreak \allowbreak 2,\allowbreak \allowbreak 3,\allowbreak \allowbreak b,\allowbreak \allowbreak e),\allowbreak \allowbreak  b\ge 5\) & 5. \\
\addlinespace[2pt]
exact crossed \((1,\allowbreak \allowbreak 2,\allowbreak \allowbreak 3,\allowbreak \allowbreak 4,\allowbreak \allowbreak e)\) & 5. \\
\addlinespace[2pt]
\(r_1=2,\allowbreak \allowbreak  j\ge 3\), canonical throughout & \(8+(\frac{64}{5})(j-3)\). \\
\end{longtable}
\endgroup


The remaining five-arm class is \((1,\allowbreak \allowbreak 2,\allowbreak \allowbreak c,\allowbreak \allowbreak b,\allowbreak \allowbreak e),\allowbreak \allowbreak  c\ge 4\). Its complete low domain \(d\le \frac{1}{64}\) has fourth loss greater than five for every \(m,\allowbreak \allowbreak L\). On \(d\ge \frac{1}{64}\) the same actual terminal regroups as

\fitdisplay{
F=(q f_2h_b)(h_cf_e),\qquad b\ge5.
}{R48.5.27}

The actual full finite continuous pair tail in \eqref{eq:r40:remaining-pair-cap} has size at most \((\frac{39}{125})\sqrt{d}\). Its scalar loss is \(16d/5\), and \(\sqrt{d}\le 8d\) here. Consequently

\fitdisplay{
h_cf_e\le\exp[-(\frac{16}{5}-\frac{312}{125})d]\le e^{-d/2}.
}{R48.5.28}

The base triple in (R48.5.27) has exactly the same receiver and connector choice as \mbox{\hyperref[part:R48.5.5]{part R48.5.5}}. Positive propagation, including both outer \(T\) edges, therefore multiplies the actual \(U\) by at most \(e^{-d/2}\). Since \(I_q(U)^3\) is homogeneous of degree eight, its bound gains \(e^{-4d}\). This bound covers the entire high domain for the fixed direction.

Every unit-minimum odd terminal has now been covered: \(r_1\ge 4,\allowbreak \allowbreak  r_1=3\), or \(r_1=2\); the last case has \(j=1,\allowbreak \allowbreak  j=2\) or \(j\ge 3\), with \(j=2\) exhausted by the three explicit classes. The analytic boundaries \(\frac{1}{64},\allowbreak \allowbreak \frac{1}{8},\allowbreak \allowbreak \frac{1}{5},\allowbreak \allowbreak \frac{1}{4}\) are included on both sides of their estimates.

\paragraph{R48.6 Every odd terminal of nonunit minimum}

We now prove the stronger same-direction statement

\fitdisplay{
\boxed{EH^4\le e^{-4d}}
 \quad\text{for every O direction in (R48.2.1)-(R48.2.2).}
}{R48.6.1}

Use the full finite packet in parts R30.2–R30.8, and A014, Sections 3–7, to prove (R48.6.1).

\paragraph{R48.6.1 Second-shortest terminal length at least four}

R2-C030 gives \(F\le \bar f_\ell \exp(R-\Gamma),\allowbreak \allowbreak  \Gamma\ge 16j d/5\). Its two actual first-edge signs have the same valid \(D\le 5d/6\), while the one finite terminal mask satisfies \(\lVert R\rVert \le \sqrt{d}/4\) and \(\sum_{i\ge 1}(D_i)_+<2d/5\). The matched integer choice \(m\) versus \(\ell+L\) makes the bare pair \(W_h \bar f_H\) with \(3\le h\le H\). Thus for every finite \(L\) the complete scalar/vector bounds are

\fitdisplay{
M\le-\frac{19}{16}+\frac{37}{15}+\frac{1}{4},
 \qquad N\le3/4.
}{R48.6.2}

Here \(C=2D+2 \sum(D_i)_+\le 37d/15\) bounds each finite propagated positive sub-sum, not an infinite actual connector. Applying (R48.1.9) to this same combined vector gives

\fitdisplay{
\text{fourth loss}\ge
 \frac{4121}{960}+(\frac{64}{5})(j-1)>4.
}{R48.6.3}

The numerator 4121 follows from \(\frac{64}{5}-\frac{367}{60}-\frac{153}{64}=\frac{4121}{960}\). This covers every \(\ell,\allowbreak \allowbreak  m,\allowbreak \allowbreak  L\) and \(j\) in this entire structural class without an additional \(d\) split.

\paragraph{R48.6.2 Second-shortest length three: the complete low packet}

Now \(\ell=2,\allowbreak \allowbreak  r_1=3\) and \(c=t_{1}\ge 4\). Put \(\eta=\beta_{c-1}\) and

\fitdisplay{
R_0=-\eta u-B^3G_{c-3}(-B)u,
 \quad\Gamma_0=d+\beta_c+\eta\ge35d/16.
}{R48.6.4}

For \(j=1,\allowbreak \allowbreak  F=\bar f_{2} f_{3} h_c\le \bar f_{2} \exp(R_0-\Gamma_{0})\). On \(d\le \frac{1}{8}\) the actual \(\bar f_{2}-\) anchored first row has coefficient one on \(BR0\), and its supplied whole signed \(\delta\) allows \(D=3d/5\). The full packet is

\fitdisplay{
\begin{gathered}
\|R_0\|\le(\frac{9}{20})\sqrt d,\\
\quad(D_1)_+\le d^2,\quad(D_2)_+\le \delta_3,\\
\quad\sum_{i\ge1}(D_i)_+<3d/7.
\end{gathered}
}{R48.6.5}

This \(\bar f_{2}\) first row is not the \(q\) first row of \mbox{\hyperref[part:R48.5]{part R48.5}}. Its actual \(\delta\) includes the additional mean \(E[(Bu)R_0]\) and the original \(\sqrt{(d+\delta_2) E R_0^2}\) cost, as explicitly paid in A014, Section 6.

The exact short \(-L\) constants are


\begingroup\small
\setlength{\tabcolsep}{3pt}
\setlength{\LTleft}{0pt}\setlength{\LTright}{\fill}
\begin{longtable}{@{}p{0.2200\linewidth}p{0.3500\linewidth}p{0.3500\linewidth}@{}}
\caption{Every longer singleton at two outer edges: parameter values.}\label{tab:parameters:r2-c048:6}\\
\toprule
Entry & \(C/d\) & \(\rho\) with \(\lVert B^{L-1}R_0\rVert \le \rho \sqrt{d}\) \\
\midrule\endfirsthead
\toprule
Entry & \(C/d\) & \(\rho\) with \(\lVert B^{L-1}R_0\rVert \le \rho \sqrt{d}\) \\
\midrule\endhead
\bottomrule\endfoot
1 & \(\frac{3}{5}\); & \(\frac{9}{20}\). \\
\addlinespace[2pt]
2 & \(\frac{53}{40}\); & \(\frac{3}{10}\). \\
\addlinespace[2pt]
3 & \(\frac{97}{60}\); & \(\frac{3}{16}\). \\
\addlinespace[2pt]
every \(L\ge 4\) & \(\frac{72}{35}\); & \(\frac{3}{25}\). \\
\end{longtable}
\endgroup


For \(L=1,\allowbreak \allowbreak L=2,\allowbreak \allowbreak L=3\) the matched bare estimate gives \(M\le -\frac{19}{16}+C/d+\rho\) and \(N\le \frac{1}{2}+\rho\). Substitution at exponent four gives respectively

\fitdisplay{
\frac{8743}{1600}>4,\qquad107/25>4,
 \qquad377/60-\frac{2057}{1024}>4.
}{R48.6.6}

The last comparison even follows from \(\frac{2057}{1024}<\frac{21}{10}\).

For every \(L\ge 4\) and \(m\ge 4\), the matched smaller index is at least four. Use the actual stronger bare pair instead:

\fitdisplay{
\begin{gathered}
M\le-\frac{219}{70}+\frac{72}{35}+\frac{3}{25}=-\frac{333}{350}<0,\\
\qquad N\le1/3+\frac{3}{25}=\frac{34}{75}<1/2.
\end{gathered}
}{R48.6.7}

The fourth loss is then at least \(\frac{35}{4}-\frac{17}{16}=\frac{123}{16}>4\). Only \(m=3\) and all \(L\ge 4\) remain in the \(j=1\) low domain. They are handled by the complete finite coefficient argument next.

\paragraph{R48.6.3 \(m=3\) and every \(L\ge 4\): retain the last two finite coefficients}

The structural rule (R48.2.2) selects \(W_{3} T^L\) because \(n=L+2>3\). Write \(G_{c-3}(-z)=U-zV\). \(At c=4\) the genuine coordinates are \(U_{1}=1,\allowbreak \allowbreak V_{1}=0\); at every other finite \(c,\allowbreak \allowbreak  0\le U,\allowbreak \allowbreak V\le 1\). The four signed moment classes are

\fitdisplay{
\begin{array}{ll}
 D_{4a}=-\int x^{2a+1}U\,d\nu,&a\ge0,\\
 D_{4a+1}=\eta \delta_{2a+1}-\int x^{2a+2}V\,d\nu,&a\ge0,\\
 D_{4a+2}=\int x^{2a+2}U\,d\nu,&a\ge0,\\
 D_{4a+3}=-\eta \delta_{2a+2}+\int x^{2a+3}V\,d\nu,&a\ge0.
 \end{array}
}{R48.6.8}

Thus every \(i\ge 2\) satisfies

\fitdisplay{
\begin{gathered}
(D_i)_+\le \delta_3,\qquad\\
\lvert D_i\rvert \le \delta_3+\eta \delta_2\\
\le(\frac{221}{1024})d<d/4,
\end{gathered}
}{R48.6.9}

where \(d\le \frac{1}{8}\) was used. These are bounds for all four residue classes, not checks of finitely many connector lengths.

The actual finite first-row propagation is

\fitdisplay{
\begin{gathered}
e^{\Gamma_0}T^LF\le Y_L\\
=\bar f_{L+2}+B^LR_0+D \bar f_{L-1}\\
+\sum_{i=1}^{L-1}D_i \bar f_{L-1-i}.
\end{gathered}
}{R48.6.10}

Combine it with the complete identity (R48.1.4) at \(L+2\). Define

\fitdisplay{
\Lambda_3=d-D,\qquad
 \Lambda_r=\beta_r-D_{r-3}\quad(4\le r\le L+2).
}{R48.L1331}

Then (R48.6.10) becomes exactly

\fitdisplay{
Y_L=1+G_{L+2}(B)u+B^LR_0
       -\sum_{r=3}^{L+2}\Lambda_r \bar f_{L+2-r}.
}{R48.6.11}

In particular the actual scalar D is absorbed only into \(\Lambda_{3}\); it has not been confused with \(D_3\). By (R48.6.5),(R48.6.9),

\fitdisplay{
\begin{gathered}
\Lambda_3=2d/5,\quad\Lambda_4\ge7d/8,\\
\quad\Lambda_r\ge(\kappa-\theta^2)d>2d/5\ (r\ge5),\\
\qquad\Lambda_r\le5d/4\ (r\ge5).
\end{gathered}
}{R48.6.12}

Every coefficient is nonnegative. We may therefore retain the last two actual terms \(\Lambda_{L+2} \bar f_{0}\) and \(\Lambda_{L+1} \bar f_1\) and discard the other nonnegative subtractions. With the exact \(W_{3}=1-d-G_{3} u\), the resulting sum \(W_{3}+Y_L\) is at most \(2+c+BV\), with

\fitdisplay{
\begin{gathered}
c=-d-\Lambda_{L+1}-\Lambda_{L+2},\\
\quad V=B^2G_{L-1}(B)u+B^{L-1}R_0-\Lambda_{L+1}1.
\end{gathered}
}{R48.6.13}

The two future coefficients retain both \(L+1\) losses. The finite odd-power mean \(E[B^3 G_{L-1}u]\) is at most \(\delta_2\). Since \(L\ge 4,\allowbreak \allowbreak  (\) R48.6.5) gives \(\lVert B^{L-1}R_0\rVert \le 3 \sqrt{d}/25\) and \(\lvert D_L\rvert \le 3d/25\). Hence

\fitdisplay{
(c+EBV)/d\le-\frac{9}{5}+\theta+\frac{3}{25}<-5/4.
}{R48.6.14}

The complete norm, including the constant, has the valid triangle bound

\fitdisplay{
\begin{gathered}
\|V\|/\sqrt d\\
\le\theta+\frac{3}{25}+(\frac{5}{4})\sqrt d\\
<\frac{5}{12}+\frac{3}{25}+\frac{15}{32}=\frac{2413}{2400}<21/20.
\end{gathered}
}{R48.6.15}

Here \(\sqrt{d}<\frac{3}{8}\) on the actual low domain. The fourth loss is at least

\fitdisplay{
\frac{35}{4}+5-(\frac{17}{4})(\frac{21}{20})^2=\frac{14503}{1600}>4.
}{R48.6.16}

This proves the entire previously weak long \(m=3\) row while retaining the original finite signed propagation and all original directions.

\paragraph{R48.6.4 All later pairs and the entire high domain}

For \(j\ge 2\) and \(d\le \frac{1}{8}\), the complete actual R2-C030/A014 \(j=2\) table has ordinary-eighth loss greater than nine for every \(L\), and each later terminal pair only improves it. Ordinary Cauchy on the same \(H\) gives ordinary-fourth loss at least \(\frac{9}{2}>4\). No weighted moment is involved.

For \(d\ge \frac{1}{8}\), use the actual centered-pair inequalities \eqref{eq:R37:centered-pair}. The complete later terminal mask \(S\) has \(\lvert S\rvert \le \sqrt{d}/4\le 3d/4; S=0\) when \(j=1\). Therefore, uniformly for every \(j\ge 1\),

\fitdisplay{
F\le \bar f_2\exp[-(\frac{1}{4}+(\frac{16}{5})(j-1))d].
}{R48.6.17}

The original selected connector maps \(\bar f_{2}\) to \(\bar f_{L+2}\) or \(W_{L+2}\). Multiplication by the selected original arm gives exactly \(W_h \bar f_H,\allowbreak \allowbreak  h=\min(m,\allowbreak \allowbreak L+2)\ge 3\). Its same-direction fourth loss is at least \(\frac{59}{16}\), as follows immediately from the actual bare sum used in R2-C030 and (R48.1.8). Thus (R48.6.17) gives total fourth loss at least

\fitdisplay{
\frac{59}{16}+1+(\frac{64}{5})(j-1)\ge75/16>4.
}{R48.6.18}

Equations (R48.6.3),(R48.6.6)-(R48.6.7),(R48.6.16),(R48.6.18) and the complete \(j\ge 2\) low table prove (R48.6.1). The analytic boundary \(d=\frac{1}{8}\) is shared, and the integer comparison \(m\) versus \(\ell+L\) is unchanged throughout.

\paragraph{R48.7 Every even terminal}

\paragraph{R48.7.1 Exact pairs and nonunit even multipairs}

Use the theorem bodies and exact direction scopes of R2-C039/A021 and R2-C041/A022. They give \(EH^4\le \exp(-6d)\) for all \(m\ge 3\) in the \(E\) and \(S\) rows of \mbox{\hyperref[part:R48.2.1]{part R48.2.1}}. \(E\) includes a single pair of minimum at least four and every larger even list of that minimum. \(S\) includes every even list of at least four arms of minimum two or three. Their already complete \(M=2\) output is retained with its exact direction; the proof uses those stated estimates.

For exact \((1,\allowbreak \allowbreak b)\), C087's actual fourth catalogue at outer length two has rate at least \(\frac{3487}{640}\), including its prescribed \(m=3,\allowbreak \allowbreak L=2\) turn. For exact \((2,\allowbreak \allowbreak b)\), the complete C043/A026 gives same-direction fourth rate \(\frac{326249}{86400}\) at every \(m\ge 3\) and \(L\ge 1\); for exact \((3,\allowbreak \allowbreak b)\), the corresponding rate is \(\frac{1901}{300}\). All these rates exceed \(\frac{29}{8}\). In particular the true short 23 and its finite \(G_{1}/G_{2}\) boundaries are included in the exact \(-(2,\allowbreak \allowbreak b)\) theorem, with the actual table in \mbox{\hyperref[part:R48.2]{part R48.2}}.

\paragraph{R48.7.2 Unit-minimum even multipairs: the actual full tail}

Let \(j\ge 2\) and

\fitdisplay{
\begin{gathered}
F=q f_{b_1}\prod_{i=2}^j h_{a_i}f_{b_i},\\
\qquad1<b_1<a_2<b_2<\cdots<a_j<b_j.
\end{gathered}
}{R48.7.1}

The complete C043/A026 finite mask proof has the following exhaustive classes and analytic domains:


\begingroup\small
\setlength{\tabcolsep}{3pt}
\setlength{\LTleft}{0pt}\setlength{\LTright}{\fill}
\begin{longtable}{@{}p{0.1200\linewidth}p{0.2667\linewidth}p{0.2667\linewidth}p{0.2667\linewidth}@{}}
\caption{Every longer singleton at two outer edges: parameter values.}\label{tab:parameters:r2-c048:7}\\
\toprule
Entry & Original integers & Low domain & High-domain full tail bound \\
\midrule\endfirsthead
\toprule
Entry & Original integers & Low domain & High-domain full tail bound \\
\midrule\endhead
\bottomrule\endfoot
I & \(b_{1}\ge 3\); & \(d\le \frac{1}{64}\); & \(\exp[-(\frac{1}{5}+(\frac{16}{5})(j-2))d]\). \\
\addlinespace[2pt]
II & \(b_{1}=2,\allowbreak \allowbreak a_{2}\ge 4\); & \(d\le \frac{1}{64}\); & same bound. \\
\addlinespace[2pt]
III & \(b_{1}=2,\allowbreak \allowbreak a_{2}=3\); & \(d\le \frac{1}{6}\); & \(\exp[-(\frac{1}{7}+(\frac{16}{5})(j-2))d]\). \\
\end{longtable}
\endgroup


In every low-domain row its actual ordinary fourth is at most \(\exp(-11d/2)\), for every \(m,\allowbreak \allowbreak L\) and every finite \(j\). These fourth-moment bounds include all signed constants.

Each high-domain tail bound in the table concerns the entire original product after the first pair. It is a pointwise cap for one finite mask. It survives complete reversal. For \(L=1\) the exact comparator is \(f_m T(p\bar f_{b_1})\), because the actual terminal in \mbox{\hyperref[part:R48.2]{part R48.2}} is \(\overline F\). For \(L\ge 2\) it is exactly \(W_m T^L(q f_{b_{1}})\). The selected directions thus coincide with the actual \(U\) directions in C087, except for \(m=3,\allowbreak \allowbreak L=2\), which is handled by \mbox{\hyperref[part:R48.3]{part R48.3}} above rather than by a turn.

Outside \(m=3,\allowbreak \allowbreak L=2\) the \(U\) fourth loss is at least \(\frac{3487}{640}\). \(At m=3,\allowbreak \allowbreak L=2\) with \(b_{1}\ge 3,\allowbreak \allowbreak  (\) R48.3.15) gives at least four. \(At m=3,\allowbreak \allowbreak L=2\) with \(b_{1}=2,\allowbreak \allowbreak  (\) R48.3.5) gives at least \(\frac{25}{8}\), with an additional negative \(4d^2\). The weakest fourth loss after multiplying the true tail is therefore

\fitdisplay{
\frac{25}{8}+\frac{4}{7}=207/56.
}{R48.7.2}

Class II instead gives \(\frac{25}{8}+\frac{4}{5}\); class I gives at least \(4+\frac{4}{5}\). Every later pair improves the fourth loss by \(\frac{64}{5}\). Including the stronger low-domain table, the entire actual unit-minimum even family satisfies

\fitdisplay{
\boxed{EH^4\le e^{-207d/56}}.
}{R48.7.3}

In particular \(2(\frac{207}{56})-\frac{36}{5}=\frac{27}{140}>\frac{1}{20}\). The same-direction all \(-T m=3,\allowbreak \allowbreak L=2\) exception is now fully paid for all terminal lengths, all finite pair counts and the whole continuous host domain.

\paragraph{R48.8 Exhaustion, uniform rate and absence of a host-selected direction}

All strict terminal lists of size at least two are even or odd. \(At\) even cardinality two, the minimum is one, two, three or at least four. \(At\) larger even cardinality, its minimum is one, two, three or at least four. Odd lists have minimum one or at least two; the corresponding subtables are given in Sections 5 and 6.

The complete minimum payment ledger is:

\begin{description}[style=nextline,leftmargin=1.6em,before=\raggedright]
\item[{\(E\) and \(S\), all \(m\ge 3\)}] Actual proved payment sufficient for \(M=2\): same-direction fourth loss 6.
\item[{Exact \((1,\allowbreak \allowbreak b)\)}] Actual proved payment sufficient for \(M=2\): C087 actual short-outer fourth loss at least \(\frac{3487}{640}\).
\item[{Exact \((2,\allowbreak \allowbreak b),\allowbreak \allowbreak  (3,\allowbreak \allowbreak b)\)}] Actual proved payment sufficient for \(M=2\): fourth losses \(\frac{326249}{86400}\) and \(\frac{1901}{300}\).
\item[{Unit-minimum even multipairs}] Actual proved payment sufficient for \(M=2\): fourth loss \(\frac{207}{56}\).
\item[{All nonunit odd O}] Actual proved payment sufficient for \(M=2\): fourth loss 4.
\item[{Unit-minimum odd first pair at least four}] Actual proved payment sufficient for \(M=2\): fourth loss at least 5.
\item[{Unit-minimum odd first pair three, one pair}] Actual proved payment sufficient for \(M=2\): fourth loss \(\frac{39}{10}\).
\item[{Unit-minimum odd first pair three, multiple pairs}] Actual proved payment sufficient for \(M=2\): fourth loss at least 7.
\item[{Exact 123}] Actual proved payment sufficient for \(M=2\): fourth loss \(\frac{29}{8}\).
\item[{Exact \(124,\allowbreak \allowbreak  L=1\)}] Actual proved payment sufficient for \(M=2\): fourth loss \(\frac{39}{10}\).
\item[{Exact 124, low nonunit \(L\)}] Actual proved payment sufficient for \(M=2\): fourth at least \(\frac{29}{8}\), or the explicit \(\frac{199}{350}\) curved output loss.
\item[{Exact \(12c,\allowbreak \allowbreak c\ge 5\), low/middle \(L=1\) and low nonunit \(L\)}] Actual proved payment sufficient for \(M=2\): fourth greater than \(\frac{29}{8}\).
\item[{Exact \(12c,\allowbreak \allowbreak c\ge 5\), high \(L=1\)}] Actual proved payment sufficient for \(M=2\): actual output loss at least \(\frac{1}{2}\).
\item[{Exact 124 or \(12c\), high nonunit \(L\)}] Actual proved payment sufficient for \(M=2\): \mbox{\hyperref[part:R48.4.4]{part R48.4.4}}, including exact curved \(W_{3}\) and opposite-first-edge \(L=3\).
\item[{Five-arm class \((1,\allowbreak \allowbreak 2,\allowbreak \allowbreak c,\allowbreak \allowbreak b,\allowbreak \allowbreak e),\allowbreak \allowbreak c\ge 4\)}] Actual proved payment sufficient for \(M=2\): low fourth \(>5\);\par\noindent high same-direction triple output plus \(4d\).
\item[{Other five-arm unit-minimum odd classes}] Actual proved payment sufficient for \(M=2\): fourth loss at least 5.
\item[{Unit-minimum odd, \(r_1=2,\allowbreak \allowbreak j\ge 3\)}] Actual proved payment sufficient for \(M=2\): fourth loss at least 8.
\end{description}

Every fourth rate at least \(\frac{29}{8}\) pays the uniform output loss \(\frac{1}{20}\) by (R48.1.10)-(R48.1.13). The explicitly curved rows have larger output losses: \(\frac{199}{350},\allowbreak \allowbreak  \frac{91}{120}\) or \(\frac{33}{20}\) in the displayed domains. The high \(L=1\) direct estimate is stronger as well. Thus the common coefficient \(\frac{1}{20}\) applies to every entry and proves Theorem \(T20\).

No threshold in \(d\) appears in \mbox{\hyperref[part:R48.2]{part R48.2}}. Each analytic split estimates the same original message on both sides and includes its endpoints. In particular no C087 turn is substituted into the C043 \(m=3,\allowbreak \allowbreak L=2\) row; no opposite-first-edge message is renamed as a forward \(U\) message; no O comparison \(\ell+L\) is replaced by \(L\) or \(L+1\); and the unit-minimum odd crossing is used only in its exact five-arm class.

